% Part: many-valued-logic
% Chapter: syntax-and-semantics
% Section: sublogics

\documentclass[../../../include/open-logic-section]{subfiles}

\begin{document}

\olfileid[id]{mvl}{syn}{sub}

\olsection{Logika Bernilai Banyak sebagai Sublogika dari~$\LogCL$}

Logika bernilai banyak yang biasa digunakan semuanya didefinisikan dengan
matriks yang nilai fungsi kebenarannya, bagi argumen dalam
$\{\True,\False\}$, sama dengan nilai fungsi kebenaran klasik. Secara khusus,
dalam logika-logika ini, jika $x\in\{\True,\False\}$, maka
$\tf{\lnot}[\Log L](x)=\tf{\lnot}[\LogCL](x)$; dan jika $\star$ merupakan
salah satu dari $\land$, $\lor$, $\lif$ serta
$x,y\in\{\True,\False\}$, maka
$\tf{\star}[\Log L](x,y)=\tf{\star}[\LogCL](x,y)$. Dengan kata lain,
fungsi kebenaran bagi $\lnot$, $\land$, $\lor$, dan~$\lif$, ketika dibatasi
pada $\{\True,\False\}$, persis sama dengan fungsi kebenaran klasik.

\begin{prop}\ollabel{prop:mvl-cl}
  Andaikan suatu logika bernilai banyak~$\Log L$ memuat konektif $\lnot$,
  $\land$, $\lor$, $\lif$ dalam bahasanya, $\True,\False\in V$, dan fungsi
  kebenarannya memenuhi:
  \begin{enumerate}
    \item\ollabel{prop:not} $\tf{\lnot}[\Log L](x)=\tf{\lnot}[\LogCL](x)$
    jika $x=\True$ atau $x=\False$;
    \item\ollabel{prop:land} $\tf{\land}[\Log L](x,y)=\tf{\land}[\LogCL](x,y)$,
    \item\ollabel{prop:lor} $\tf{\lor}[\Log L](x,y)=\tf{\lor}[\LogCL](x,y)$,
    \item\ollabel{prop:lif} $\tf{\lif}[\Log L](x,y)=\tf{\lif}[\LogCL](x,y)$,
    jika $x,y\in\{\True,\False\}$.
  \end{enumerate}
  Maka, bagi setiap valuasi~$\pAssign v$ ke dalam~$V$ sedemikian sehingga
  $\pAssign v(p)\in\{\True,\False\}$, dan bagi setiap formula dalam
  fragmen klasik yang dibangun dari variabel proposisional hanya dengan
  keempat konektif dalam hipotesis, berlaku
  $\pValue v[\Log L](!A)=\pValue v[\LogCL](!A)$.
\end{prop}

\begin{proof}
  Dengan induksi pada~$!A$.
  \begin{enumerate}
  \item Jika $!A\ident p$ bersifat atomik, kita mempunyai
  $\pValue v[\Log L](!A)=\pAssign v(p)=\pValue v[\LogCL](!A)$.
  \item Jika $!A\ident\lnot B$, kita mempunyai
  \begin{align*}
    \pValue v[\Log L](!A) & = \tf{\lnot}[\Log L](\pValue v[\Log L](!B)) & &\text{berdasarkan \olref[val]{defn:pValue}}\\
    & = \tf{\lnot}[\Log L](\pValue v[\LogCL](!B)) && \text{berdasarkan hipotesis induksi}\\
    & = \tf{\lnot}[\LogCL](\pValue v[\LogCL](!B))
&&\text{berdasarkan asumsi \olref{prop:not},}\\
&&&\text{karena $\pValue v[\LogCL](!B)\in\{\True,\False\}$,}\\
& = \pValue v[\LogCL](!A) &&\text{berdasarkan \olref[val]{defn:pValue}}.
\end{align*}
\item Jika $!A\ident(!B\land !C)$, kita mempunyai
\begin{align*}
  \pValue v[\Log L](!A) & = \tf{\land}[\Log L](\pValue v[\Log L](!B), \pValue v[\Log L](!C)) & &\text{berdasarkan \olref[val]{defn:pValue}}\\
  & = \tf{\land}[\Log L](\pValue v[\LogCL](!B),\pValue v[\LogCL](!C)) && \text{berdasarkan hipotesis induksi}\\
  & = \tf{\land}[\LogCL](\pValue v[\LogCL](!B),\pValue v[\LogCL](!C))
&&\text{berdasarkan asumsi \olref{prop:land},}\\
&&&\text{karena $\pValue v[\LogCL](!B),\pValue v[\LogCL](!C)\in\{\True,\False\}$,}\\
& = \pValue v[\LogCL](!A) &&\text{berdasarkan \olref[val]{defn:pValue}}.
\end{align*}
\end{enumerate}
Kasus $!A\ident(!B\lor !C)$ dan $!A\ident(!B\lif !C)$ serupa.
\end{proof}

\begin{cor}
  Jika suatu logika bernilai banyak memenuhi syarat-syarat dalam
  \olref{prop:mvl-cl}, $\True\in V^+$, dan $\False\notin V^+$, maka pada
  fragmen klasik bersama berlaku
  ${\Entails[\Log L]}\subseteq{\Entails[\LogCL]}$; yakni, jika
  $\Gamma\Entails[\Log L]!B$, maka $\Gamma\Entails[\LogCL]!B$. Secara
  khusus, setiap tautologi dari~$\Log L$ dalam fragmen tersebut juga
  merupakan tautologi klasik.
\end{cor}

\begin{proof}
  Kita membuktikan kontraposisinya. Andaikan
  $\Gamma\Entails/[\LogCL]!B$. Maka terdapat suatu
  !!{valuation}~$\pAssign v\colon\PVar\to\{\True,\False\}$ sedemikian
  sehingga $\pValue v[\LogCL](!A)=\True$ bagi setiap $!A\in\Gamma$ dan
  $\pValue v[\LogCL](!B)=\False$. Karena $\True,\False\in V$,
  !!{valuation}~$\pAssign v$ juga merupakan !!a{valuation} bagi~$\Log L$.
  Berdasarkan \olref{prop:mvl-cl}, $\pValue v[\Log L](!A)=\True$ bagi setiap
  $!A\in\Gamma$ dan $\pValue v[\Log L](!B)=\False$. Karena
  $\True\in V^+$ dan $\False\notin V^+$, berarti
  $\pSat{v}{\Gamma}[\Log L]$ dan $\pSat/{v}{!B}[\Log L]$; dengan kata lain,
  $\Gamma\Entails/[\Log L]!B$.
\end{proof}
\end{document}
